\chapter{自旋与角动量}

\textbf{6-1}\quad 电子偶素 $(e^{+}e^{-}$ 束缚态)类似于氢原子, 只是用一个正电子代替质子作为核. 在非相对论近似下, 其能量和波函数同氢原子相似. 今设在电子偶素的基态 (s态)里, 存在一种接触型自旋交换作用 $\hat{H}' = -\dfrac{8\pi}{3}\hat{M}_{\mathrm{p}} \cdot \hat{M}_{\mathrm{e}}\delta^{3}(\boldsymbol{r})$ , 其中 $\hat{M}_{\mathrm{p}} = \dfrac{e}{mc}\hat{S}_{\mathrm{p}}$ 与 $\hat{M}_{\mathrm{e}} = -\dfrac{e}{mc}\hat{S}_{\mathrm{e}}$ 分别是正负电子的自旋磁矩. 利用一级微扰论计算此基态中自旋单态与三重态之间的能量差, 决定哪一能量更低. 已知氢原子基态波函数

$$
\psi_ {1 0 0} (\boldsymbol {r}) = \sqrt {\frac {1}{\pi a ^ {3}}} \mathrm{e} ^ {- r / a}, a = \frac {\hbar^ {2}}{\mu e ^ {2}}, \frac {e ^ {2}}{\hbar c} = \frac {1}{1 3 7}
$$

\vfill

\newpage

\textbf{6-2}\quad 电子偶素是由正负电子构成的类氢原子体系. 考虑处于基态的电子偶素 $(l = 0)$ , 系统的哈密顿量可写成 $\hat{H} = \hat{H}_0 + \hat{H}_s$ , 其中 $\hat{H}_0$ 是通常的与自旋无关的库仑力部分, $\hat{H}_s = A\hat{S}_p \cdot \hat{S}_e$ 是正电子与负电子的自旋作用部分. 请问在无外磁场作用时, 选择怎样的自旋和角动量本征态最方便? 对这些态, 计算由于 $\hat{H}_s$ 引起的能量的改变.

\vfill

\newpage

\textbf{6-3}\quad 均匀磁场中电子偶素(电子-正电子束缚态)的哈密顿量(不考虑空间运动)为 $\hat{H} = J\sigma_{1} \cdot \sigma_{2} + FB(\sigma_{1z} - \sigma_{2z})$ ，其中 $J$ 与 $F$ 是常实数， $B$ 是磁场强度，下标1与2分别代表电子与正电子， $\sigma$ 是泡利矩阵。(1) 当 $B = 0$ 时， $\hat{H}$ 的本征函数与本征值是什么？(2) 当 $B \neq 0$ 时， $\hat{H}$ 的本征函数与本征值是什么？

\vfill

\newpage

\textbf{6-4}\quad 一束极化的 s 波 (l=0) 中子通过一个不均匀的磁场后分裂成强度不同的两束, 其中自旋反平行于磁场的一束与自旋平行于磁场的一束的强度之比为 3:1. 求入射中子自旋方向与磁场方向夹角的大小.

\vfill

\newpage

\textbf{6-5}\quad $\sigma_{1},\sigma_{2}$ 为泡利矩阵, 证明 $\mathbf{e}^{\mathrm{i}\alpha \sigma_j} = \cos \alpha +\mathrm{i}\sigma_j\sin \alpha ,j = 1,2,\alpha$ 为实数并推广到矩阵 $\pmb {\sigma} = \lambda \pmb{\sigma}_{1} + \mu \pmb{\sigma}_{2}(\lambda^{2} + \mu^{2} = 1)$ 的情形.

\vfill

\newpage

